\chapter{Waves and optics essentials}\label{ch:waves}

Waves are what oscillations look like spread across space: chapter~7's
single bob multiplied into a string, a surface, a field. Standing waves,
interference and diffraction are all one fact --- phase adds --- read at
three scales.

\section{Wave quantities: $v=\lambda f$}\label{sec:wav-quant}

A travelling wave $y(x,t)=A\cos(kx-\omega t)$ shifts its whole shape
at speed $v=\omega/k=\lambda f$ (figure~\ref{fig:trav}). Which of
$\lambda,f,v$ changes on entering a new medium is the first question to
ask: $f$ is invariant (the source sets it), $v$ is the medium's, so
$\lambda$ adjusts.

\begin{figure}[htbp]\centering
\begin{tikzpicture}[scale=0.85,font=\scriptsize]
\foreach \sh/\lab in {0/{$t=0$},-1.9/{$t=T/4$},-3.8/{$t=T/2$}}{
\begin{scope}[yshift=\sh cm]
\draw[very thick,accent,domain=0:6,samples=60] plot (\x,{0.5*sin(360*\x/3)}) ;
\draw[->] (0,-0.8) -- (6.2,-0.8);
\node at (7.0,0) {\lab};
\end{scope}}
\node[align=left] at (3,-5.2) {$y(x,t)=A\cos(kx-\omega t)$:\\the whole shape slides right};
\end{tikzpicture}

\caption{Three snapshots of a travelling wave: the shape slides, no
particle accompanies it.}
\label{fig:trav}
\end{figure}

\section{Waves on strings and reflections}\label{sec:wav-string}

On a string $v=\sqrt{T/\mu}$: tension makes speed, mass density fights
it. At a fixed end the reflected pulse inverts; at a free end it
returns upright (figure~\ref{fig:reflect}) --- boundary conditions
already deciding which harmonics can exist.

\begin{figure}[htbp]\centering
\begin{tikzpicture}[scale=0.8,font=\scriptsize]
\draw[very thick] (5.4,-0.9) -- (5.4,0.9);
\draw[accent,thick] (0,0) .. controls (2,0.6) and (3.5,0.9) .. (5.4,0);
\draw[amber,thick,dashed] (0,0) .. controls (2,-0.6) and (3.5,-0.9) .. (5.4,0);
\node[align=left] at (2.7,-1.6) {fixed end: reflected pulse inverts};
\begin{scope}[xshift=7.5cm]
\draw[thick] (5.4,0.7) -- (5.4,-0.7);
\draw[accent,thick] (0,0) .. controls (2,0.6) and (3.5,0.9) .. (5.4,0);
\draw[amber,thick,dashed] (0,0) .. controls (2,0.6) and (3.5,0.9) .. (5.4,0);
\node[align=left] at (2.7,-1.6) {free end: reflected pulse upright};
\end{scope}
\end{tikzpicture}

\caption{Reflection rules: fixed end inverts, free end returns
upright.}
\label{fig:reflect}
\end{figure}

\section{Standing waves and harmonics}\label{sec:wav-stand}

Two counter-propagating copies make a standing wave:
$y=2A\sin kx\cos\omega t$ --- positions frozen (nodes), the rest
swinging in place (figure~\ref{fig:stand}). A pinned--pinned string of
length $L$ hosts only $\lambda_n=2L/n$, i.e.\
$f_n=nv/2L$: the spectrum of every instrument.

\begin{figure}[htbp]\centering
\begin{tikzpicture}[scale=0.9,font=\scriptsize]
\draw[very thick,accent,domain=0:6.28,samples=80] plot (\x,{0.9*sin(180*\x)}) ;
\draw[very thick,dashed,accent!50,domain=0:6.28,samples=80] plot (\x,{-0.9*sin(180*\x)}) ;
\draw[->] (0,0) -- (6.6,0);
\foreach \x/\l in {0/N,3.14/N,6.28/N}{\draw[thick] (\x,-0.15) -- (\x,0.15) node[above]{\l};}
\foreach \x/\l in {1.57/A,4.71/A}{\draw[thick] (\x,-0.15) -- (\x,0.15) node[above]{\l};}
\node[align=left] at (3.14,-1.6) {n-th harmonic: $f_n=nv/2L$,\\$\lambda_n=2L/n$};
\end{tikzpicture}

\caption{The standing-wave envelope: nodes still, antinodes at full
swing.}
\label{fig:stand}
\end{figure}

\begin{example}[Which modes a string supports]\label{ex:modes}
A $0.80$\,m string under tension gives $v=240$\,m/s. List its first
three resonant frequencies and the mode shapes.
\par\smallskip
\textbf{Strategy.} $f_n=nv/2L$.
\begin{solution}
$f_1=240/1.6=150$\,Hz, $f_2=300$\,Hz, $f_3=450$\,Hz. The $n$-th
mode fits $n$ half-wavelengths between the pins. \checkmark
\end{solution}
\end{example}

\section{Sound and the Doppler shift}\label{sec:wav-dopp}

Sound is the same mathematics in air ($v\approx343$\,m/s): source
moving compresses fronts ahead, stretches behind
(figure~\ref{fig:dopp});
$f'=f\,v/(v\mp v_s)$ for a moving source, $f'=f(v\pm v_o)/v$ for a
moving observer.

\begin{figure}[htbp]\centering
\begin{tikzpicture}[scale=0.6,font=\scriptsize]
\foreach \r in {0.5,1.0,1.5,2.0}{\draw[accent!60] (0.6,0) circle (\r);}
\fill (0.6,0) circle (0.1);
\draw[->,thick,amber] (0.6,0) -- (1.8,0);
\node[align=left] at (1.4,-2.6) {source moving right:\\compressed ahead,\\stretched behind};
\end{tikzpicture}

\caption{Doppler geometry: fronts bunch ahead, spread behind.}
\label{fig:dopp}
\end{figure}

\begin{example}[Doppler both directions]\label{ex:dopp}
An ambulance's siren is $600$\,Hz; it passes at $30$\,m/s. What does a
stationary listener hear approaching and receding ($v=343$\,m/s)?
\par\smallskip
\textbf{Strategy.} Moving source, stationary observer:
$f'=fv/(v\mp v_s)$.
\begin{solution}
Approaching: $600\times343/313=658$\,Hz; receding:
$600\times343/373=552$\,Hz. The $\sim100$\,Hz drop on pass is the
signature you hear. \checkmark
\end{solution}
\end{example}

\section{Interference: the two-slit geometry}\label{sec:wav-2slit}

Light arriving through two slits reaches a screen point with a path
difference $d\sin\theta$ (figure~\ref{fig:2slit}); bright where it is
an integer $\lambda$, dark where half-integer. Fringe spacing
$\Delta y=\lambda L/d$ --- geometry turning nanometres into
millimetres.

\begin{figure}[htbp]\centering
\begin{tikzpicture}[scale=0.9,font=\scriptsize]
\draw[very thick] (0,-1.2) -- (0,1.2);
\fill (0,-0.4) circle (0.08); \fill (0,0.4) circle (0.08);
\foreach \r in {1.2,2.4,3.6}{\draw[accent!50,dashed] (0,0.4) arc (20:-20:\r);
\draw[accent!50,dashed] (0,-0.4) arc (20:-20:\r);}
\draw[->,thick] (0,0.4) -- (4.6,0.1);
\draw[->,thick] (0,-0.4) -- (4.6,0.1);
\node[align=left] at (2.4,-1.6) {$d\sin\theta=m\lambda$:\\path difference rules the fringe};
\end{tikzpicture}

\caption{Two-slit geometry: the phase difference is geometric.}
\label{fig:2slit}
\end{figure}

\begin{example}[Two-slit fringe spacing]\label{ex:2slit}
$\lambda=633$\,nm light through $d=0.20$\,mm slits hits a screen
$2.0$\,m away: find the fringe spacing.
\par\smallskip
\textbf{Strategy.} $\Delta y=\lambda L/d$.
\begin{solution}
$\Delta y=633\times10^{-9}\times2.0/2.0\times10^{-4}=6.3$\,mm ---
comfortable to see, and measuring it back-computes $\lambda$.
\checkmark
\end{solution}
\end{example}

\begin{example}[Three slits: phasors do the counting]\label{ex:3slit}
Add a third slit at the same spacing $d$: find the bright and dark
positions on the screen and the strength of the first
\emph{secondary} maximum.
\par\smallskip
\textbf{Strategy.} Each slit contributes the same phasor length
$E_0$, rotated by $\varphi=(2\pi/\lambda)d\sin\theta$ from its
neighbour; add the three arrows head-to-tail.
\begin{solution}
$|E_{\rm tot}|^{2}=E_0^{2}\left[\dfrac{\sin(3\varphi/2)}{\sin(\varphi/2)}\right]^{2}$.
Principal maxima still at $\varphi=2\pi m$ --- i.e.\ $d\sin\theta=m\lambda$,
\emph{unchanged} from two slits --- but now two minima sit between
them, at $\varphi=\pm2\pi/3,\pm4\pi/3$
(figure~\ref{fig:phasor3}), and a weak secondary maximum stands
midway at $\varphi=\pi$ where the phasors leave one $E_0$ standing:
$1/9$ the principal intensity. For $d=0.05$\,mm, $\lambda=600$\,nm,
$L=1$\,m: principal spots at $1.2$\,cm spacing, darkness at
$0.4$ and $0.8$\,cm in between. Every added slit makes the bright
lines \emph{sharper} without moving them --- that sharpening is the
grating's whole trick. \checkmark
\end{solution}
\end{example}

\begin{figure}[htbp]\centering
\begin{tikzpicture}[scale=0.8,font=\scriptsize]
% phi = 0: collinear, maximum
\begin{scope}
\draw[->,thick,black] (0,0) -- (0.9,0);
\draw[->,thick,black] (0.9,0) -- (1.8,0);
\draw[->,thick,black] (1.8,0) -- (2.7,0);
\draw[->,very thick,accent] (0,-0.35) -- (2.7,-0.35);
\node[align=center] at (1.35,-0.75) {$\varphi=0$: sum $=3E_0$ (bright)};
\end{scope}
% phi = 2pi/3: closed triangle, minimum
\begin{scope}[xshift=4.4cm]
\draw[->,thick,black] (0,0) -- (0.9,0);
\draw[->,thick,black] (0.9,0) -- ({0.9+0.9*cos(120)},{0.9*sin(120)});
\draw[->,thick,black] ({0.9+0.9*cos(120)},{0.9*sin(120)}) -- (0,0);
\node[align=center] at (0.5,-0.75) {$\varphi=2\pi/3$: closes (dark)};
\end{scope}
% phi = pi: two cancel, one survives -> 1/9 intensity
\begin{scope}[xshift=8.2cm]
\draw[->,thick,black] (0,0) -- (0.9,0);
\draw[->,thick,black] (0.9,0) -- (0,0);
\draw[->,thick,black] (1.15,0) -- (2.05,0);
\draw[->,very thick,accent] (0,-0.35) -- (0.9,-0.35);
\node[align=center] at (1.0,-0.75) {$\varphi=\pi$: sum $=E_0$ ($\tfrac19$ light)};
\end{scope}
\end{tikzpicture}

\caption{Three-slit phasor bookkeeping: aligned arrows give the
principal maximum; the equilateral closure gives the new minima; one
surviving arrow is the $1/9$-strength secondary.}
\label{fig:phasor3}
\end{figure}

\section{Diffraction: one slit is enough}\label{sec:wav-diffr}

Even a single slit spreads the beam: the first zero sits at
$a\sin\theta=\lambda$ (figure~\ref{fig:1slit}), so narrower $a$ spreads
\emph{more}. Rayleigh's criterion: two point sources resolve when the
first's minimum sits on the second's maximum --- aperture size as the
optics budget.

\begin{figure}[htbp]\centering
\begin{tikzpicture}
\begin{axis}[width=8.2cm,height=4.6cm,xlabel={$\theta$},ylabel={$I$},
xmin=-3,xmax=3,ymin=0,ymax=1.15,font=\small,xtick=\empty,ytick=\empty]
\addplot[very thick,accent,domain=-2.95:2.95,samples=150] {(sin(deg(x))/deg(x))^2};
\end{axis}
\node[align=left] at (3.4,-0.6) {$a\sin\theta=\lambda$:\\first zero sets the width};
\end{tikzpicture}

\caption{Single-slit intensity: the first zero at
$a\sin\theta=\lambda$ sets the beam's natural width.}
\label{fig:1slit}
\end{figure}

\begin{example}[Diffraction limit estimate]\label{ex:rayleigh}
A $2.4$\,m telescope mirror at $\lambda=550$\,nm: the angular
resolution floor?
\par\smallskip
\textbf{Strategy.} Circular aperture: $\theta_{\rm min}\approx1.22\lambda/D$.
\begin{solution}
$\theta_{\rm min}=1.22(5.5\times10^{-7})/2.4=2.8\times10^{-7}$\,rad
$\approx0.06$ arc-seconds --- Hubble's diffraction floor. \checkmark
\end{solution}
\end{example}

\section{Thin films and the phase flip}\label{sec:wav-film}

A soap bubble's colour: reflection off the front face flips phase by
$\pi$ (higher index), the back face doesn't --- so bright when
$2nt=(m+\tfrac12)\lambda$: the film's thickness painted as colour.

\begin{example}[Soap-film colour]\label{ex:film}
A film $n=1.33$ shows green ($\lambda=550$\,nm) brightest in
reflection at normal incidence. Thinnest thickness?
\par\smallskip
\textbf{Strategy.} Front flip + back none: $2nt=(m+\tfrac12)\lambda$.
\begin{solution}
$2nt=\lambda/2\Rightarrow t=\lambda/4n=103$\,nm. \checkmark
\end{solution}
\end{example}

\begin{lstlisting}[language=Python,caption={\texttt{code/ch16\_superposition.py} --- standing-wave snapshots + beat envelope.},label={lst:sup}]
return math.sin(k*x-w*t) + math.sin(k*x+w*t)
\end{lstlisting}

Output of a real run:
\begin{lstlisting}[numbers=none]
standing wave snapshots y(x,t), x in [0,1]:
 t=0.000: +0.00 +1.41 +2.00 +1.41 +0.00 -1.41 -2.00 -1.41 -0.00
 t=0.125: +0.00 +1.00 +1.41 +1.00 +0.00 -1.00 -1.41 -1.00 -0.00
 t=0.250: +0.00 +0.00 +0.00 +0.00 +0.00 -0.00 -0.00 -0.00 +0.00
node check x=0.5: True
antinode x=0.25 swing: [2.0, 1.41, 0.0]
beat: f1 = 440  f2 = 444
env at t=0,.06,.125,.19,.25,.31,.375: [1.82, 1.46, 0.43, 1.45, 1.95, 1.46, 0.43]
envelope nodes at t = [0.125, 0.375, 0.625]
measured beat period = node spacing = 0.25  theory 1/|f1-f2| = 0.25
\end{lstlisting}

The snapshots freeze mid-swing exactly on the half-cycle; nodes sit
exactly still; the beat envelope's nodes land every
$1/(2\Delta f)=0.125$\,s --- the 4\,Hz warble you hear when two tuners
disagree.

\begin{example}[Headlights at night: your pupil's Rayleigh limit]\label{ex:headlight}
How far can your unaided eye ($D\approx5$\,mm pupil) resolve a car's
two headlights ($1.5$\,m apart) at $\lambda=550$\,nm?
\par\smallskip
\textbf{Strategy.} Rayleigh for a circular aperture:
$\theta_{\rm min}=1.22\lambda/D$; then $\theta\approx s/L$.
\begin{solution}
$\theta_{\rm min}=1.22\times5.5\times10^{-7}/5\times10^{-3}
=1.34\times10^{-4}$\,rad;
$L=s/\theta=1.5/1.34\times10^{-4}\approx11$\,km.
Inside that distance the pair splits into two; beyond it, one smeared
blob --- every dark road is a diffraction demonstration.
\checkmark
\end{solution}
\end{example}

\begin{lstlisting}[language=Python,caption={\texttt{code/ch16\_nslit.py} --- $N$-slit intensity: principal maxima fixed, lines sharpen, satellites counted.},label={lst:nslit}]
def intensity(phi, N):
    return (np.sin(N * phi / 2) / np.sin(phi / 2))**2
phi = np.linspace(1e-6, 2 * np.pi - 1e-6, 400001)
for N in [2, 3, 8]:
    I = intensity(phi, N)
    seg = I[1:]
    peaks = int(np.sum((seg[1:-1] > seg[:-2]) & (seg[1:-1] > seg[2:])
                       & (seg[1:-1] > 1)))
    print(f"N={N}: principal max = {N*N:3d} E0^2,"
          f"  secondary maxima measured = {peaks}")
\end{lstlisting}

Output of a real run:
\begin{lstlisting}[numbers=none]
N=2: principal max =   4 E0^2,  secondary maxima measured = 0
N=3: principal max =   9 E0^2,  secondary maxima measured = 1
N=8: principal max =  64 E0^2,  secondary maxima measured = 6
N=2: first zero at sin(theta) = lambda/(2 d)   line width ~ 1/2
N=3: first zero at sin(theta) = lambda/(3 d)   line width ~ 1/3
N=8: first zero at sin(theta) = lambda/(8 d)   line width ~ 1/8
\end{lstlisting}

The machine confirms what the phasors sketched: principal maxima
carry $N^{2}$ units of $E_0^{2}$, exactly $N-2$ secondary maxima
live between them, and the first zero pulls in like $1/N$ ---
slit count buys sharpness \emph{and} brightness together, which is
why a $500$-line grating resolves what three slits cannot.

\section{In engineering practice: optics on the bench}\label{sec:wav-eng}

Antenna arrays are two-slit geometry at radio wavelength; loudspeaker
beaming is figure~\ref{fig:1slit}'s diffraction floor; interferometers
(LIGO's arms) are \ref{ex:2slit}'s path-difference bookkeeping pushed
to $\lambda/1000$ precision; and every lens's resolving power hits
Rayleigh's aperture budget.

\section{Common mistakes}\label{sec:wav-mistakes}

\begin{tripwire}[$v=\lambda f$ misused across media]
$f$ is invariant (source), $v$ is medium's, $\lambda$ adjusts --- swap
the wrong two and every answer follows the error.
\end{tripwire}
\begin{tripwire}[Node vs antinode at a fixed end]
Fixed = node = inversion; free = antinode = upright. Mixing them
breaks the harmonic count.
\end{tripwire}
\begin{tripwire}[Interference vs diffraction blurred]
Interference = few sources adding; diffraction = one aperture's
continuum. A single slit does the second.
\end{tripwire}
\begin{tripwire}[Forgetting the thin-film phase flip]
The front-face flip makes the conditions ``backwards'': bright
reflection needs the \emph{half}-integer path.
\end{tripwire}
\begin{tripwire}[Doppler numerator/denominator swap]
Source moves $\to$ denominator; observer moves $\to$ numerator.
Deriving from figure~\ref{fig:dopp} is safer than the formula.
\end{tripwire}

\section{Summary}\label{sec:wav-summary}

\begin{table}[ht]\centering
\begin{tabularx}{\linewidth}{@{}lY@{}}
\toprule
\tablehead\thead{Idea} & \thead{One-line version}\\
\midrule
travelling & $v=\lambda f$; $f$ invariant, $\lambda$ adjusts \\
string & $v=\sqrt{T/\mu}$; fixed inverts, free upright \\
standing & $\lambda_n=2L/n$; nodes/antinodes \\
Doppler & source in denominator, observer in numerator \\
two-slit & $d\sin\theta=m\lambda$; $\Delta y=\lambda L/d$ \\
diffraction & $a\sin\theta=\lambda$ first zero; Rayleigh floor \\
\bottomrule
\end{tabularx}
\caption{Waves/optics cheat sheet.}
\label{tab:wavsummary}
\end{table}

\begin{wins}[Definition of done]
You can now: pick which quantity stays fixed across media, count
harmonics for both end conditions, compute fringe spacings and first
zeros, derive thin-film conditions with the flip, and run the
superposition/beats checks.
\end{wins}

\chapternotes{%
The order follows OpenStax vols.~1 and 3 \cite{openstax1} with
Feynman~I \S30, \S47--51 the inspiration \cite{feynman1}; MIT 8.03SC
for the full wave treatment \cite{ocw801}. The numbers are the
listing's real output.}

\begin{rulebox}[{Watch instead}]
\begin{enumerate}
\item \emph{Standing waves on strings} --- Khan Academy Physics,
13:25, \url{https://youtu.be/gT0IqL1dyyk} ---
\S\ref{sec:wav-string}--\S\ref{sec:wav-stand} (accessed 2026-10-11).
\item \emph{Young's double slit introduction} --- Khan Academy, 9:33,
\url{https://youtu.be/Pk6s2OlKzKQ} --- \S\ref{sec:wav-2slit} quickly
(accessed 2026-10-11).
\end{enumerate}
\end{rulebox}

\section*{Further reading}
\addcontentsline{toc}{section}{Further reading}

\begin{itemize}
\item MIT OCW \emph{8.03SC Vibrations and Waves}. At
\url{https://ocw.mit.edu}, course ``8.03SC'' (verified 2026-10-11).
\item OpenStax \emph{College Physics 2e}, chs.~16--17 (waves, sound)
and \emph{University Physics} vol.~3 (optics) ---
\url{https://openstax.org} (verified 2026-10-11).
\item R.~P. Feynman, \emph{The Feynman Lectures on Physics} I,
\S30, \S47--51 --- interference and diffraction at their most lyrical
(in print).
\item Khan Academy, \emph{Waves and optics} --- drills matching
\S\ref{sec:wav-quant}--\S\ref{sec:wav-2slit}: verified 2026-10-11 at
\url{https://www.khanacademy.org} (science $\to$ physics $\to$ waves).
\item PhET, \emph{Wave Interference} interactive --- slits, films and
fringes you can drag: at \url{https://phet.colorado.edu} (verified
2026-10-11).
\end{itemize}

\section*{Additional exercises}
\addcontentsline{toc}{section}{Additional exercises}

\begin{exercise}[easy]
A wave has $v=340$\,m/s, $f=680$\,Hz: find $\lambda$. (AP-C)
\end{exercise}
\begin{exercise}[easy]
Mark nodes and antinodes for the third harmonic on a pinned--pinned
string. (AP-C)
\end{exercise}
\begin{exercise}[easy]
What is the mode number of a string vibrating at $4f_1$? (AP-C)
\end{exercise}
\begin{exercise}[easy]
A source approaches you: does $f'$ go up or down? Which factor?
(F=ma)
\end{exercise}
\begin{exercise}[easy]
Two slits are $0.1$\,mm apart, screen $1$\,m: find fringe spacing for
$\lambda=500$\,nm. (AP-C)
\end{exercise}
\begin{exercise}[medium]
Find the tension that raises a $0.6$\,kg/m string's fundamental to
$200$\,Hz. (AP-C FR-style)
\end{exercise}
\begin{exercise}[medium]
Derive the condition for dark fringes in two-slit and explain why
they're perfectly dark. (AP-C)
\end{exercise}
\begin{exercise}[medium]
A soap film goes red $\to$ green $\to$ blue: what is changing and in
which direction? (APhO-lite)
\end{exercise}
\begin{exercise}[medium]
Two tuning forks beat at $3$\,Hz: find both possible $f_2$ if
$f_1=440$\,Hz. (F=ma)
\end{exercise}
\begin{exercise}[medium]
Why does narrowing a slit widen the diffraction pattern? Argue from
Huygens. (AP-C)
\end{exercise}
\begin{exercise}[hard]
Estimate the angular resolution of your eye's pupil ($D\approx5$\,mm)
and compare with the textbook resolution test. (estimation)
\end{exercise}
\begin{exercise}[easy]
In two-slit interference, what happens to the fringe spacing if the
screen is moved twice as far --- and to the fringe \emph{brightness}
if the slits are moved twice as close? (AP-C)
\end{exercise}
\begin{exercise}[medium]
A grating has $N$ slits at spacing $d$: argue the principal maxima
sharpen like $1/N$ and estimate the resolving power
$\lambda/\Delta\lambda\sim mN$ in order $m$. (USAPhO-style)
\end{exercise}
\begin{exercise}[medium]
A siren passes a stationary listener: derive the heard frequency on
approach and recession, and explain why the two shifts are not
symmetric. (F=ma)
\end{exercise}
\begin{exercise}[medium]
Two antennas radiate in phase at $3$\,MHz, $500$\,m apart: find the
interference pattern on a distant perpendicular bisector. (engineering)
\end{exercise}
\begin{exercise}[medium]
An open--open pipe vs a closed--open pipe: derive both harmonic
series. (AP-C)
\end{exercise}
\begin{exercise}[medium]
A light wave enters glass ($n=1.5$): find what stays fixed and what
changes among $f,\lambda,v$. (AP-C)
\end{exercise}
\begin{exercise}[medium]
Why can a lens's resolution never beat its diffraction floor? Argue
via aperture $D$. (concept)
\end{exercise}
